%%%%%%%%%%%%%%%%%%%%%%%%%%%%%%%%%%%%%%%%%
% Lachaise Assignment
% LaTeX Template
% Version 1.0 (26/6/2018)
%
% This template originates from:
% http://www.LaTeXTemplates.com
%
% Authors:
% Marion Lachaise & François Févotte
% Vel (vel@LaTeXTemplates.com)
%
% License:
% CC BY-NC-SA 3.0 (http://creativecommons.org/licenses/by-nc-sa/3.0/)
% 
%%%%%%%%%%%%%%%%%%%%%%%%%%%%%%%%%%%%%%%%%

%----------------------------------------------------------------------------------------
%	PACKAGES AND OTHER DOCUMENT CONFIGURATIONS
%----------------------------------------------------------------------------------------

\documentclass{article}
\usepackage{amsthm}
\usepackage{listings}
\usepackage{xcolor}
\usepackage{graphicx}


\lstset{
	language=Python,
	basicstyle=\ttfamily\small,
	keywordstyle=\color{blue},
	stringstyle=\color{red},
	commentstyle=\color{purple},
	showstringspaces=false,
	numbers=left,
	numberstyle=\tiny\color{gray},
	breaklines=true,
	frame=single,
	captionpos=b
}
\newtheorem{definition}{Definition}[section]

%%%%%%%%%%%%%%%%%%%%%%%%%%%%%%%%%%%%%%%%%
% Lachaise Assignment
% Structure Specification File
% Version 1.0 (26/6/2018)
%
% This template originates from:
% http://www.LaTeXTemplates.com
%
% Authors:
% Marion Lachaise & François Févotte
% Vel (vel@LaTeXTemplates.com)
%
% License:
% CC BY-NC-SA 3.0 (http://creativecommons.org/licenses/by-nc-sa/3.0/)
% 
%%%%%%%%%%%%%%%%%%%%%%%%%%%%%%%%%%%%%%%%%

%----------------------------------------------------------------------------------------
%	PACKAGES AND OTHER DOCUMENT CONFIGURATIONS
%----------------------------------------------------------------------------------------

\usepackage{amsmath,amsfonts,stmaryrd,amssymb} % Math packages

\usepackage{enumerate} % Custom item numbers for enumerations

\usepackage[ruled]{algorithm2e} % Algorithms

\usepackage[framemethod=tikz]{mdframed} % Allows defining custom boxed/framed environments

\usepackage{listings} % File listings, with syntax highlighting
\lstset{
	basicstyle=\ttfamily, % Typeset listings in monospace font
}

%----------------------------------------------------------------------------------------
%	DOCUMENT MARGINS
%----------------------------------------------------------------------------------------

\usepackage{geometry} % Required for adjusting page dimensions and margins

\geometry{
	paper=a4paper, % Paper size, change to letterpaper for US letter size
	top=2.5cm, % Top margin
	bottom=3cm, % Bottom margin
	left=2.5cm, % Left margin
	right=2.5cm, % Right margin
	headheight=14pt, % Header height
	footskip=1.5cm, % Space from the bottom margin to the baseline of the footer
	headsep=1.2cm, % Space from the top margin to the baseline of the header
	%showframe, % Uncomment to show how the type block is set on the page
}

%----------------------------------------------------------------------------------------
%	FONTS
%----------------------------------------------------------------------------------------

\usepackage[utf8]{inputenc} % Required for inputting international characters
\usepackage[T1]{fontenc} % Output font encoding for international characters

\usepackage{XCharter} % Use the XCharter fonts

%----------------------------------------------------------------------------------------
%	COMMAND LINE ENVIRONMENT
%----------------------------------------------------------------------------------------

% Usage:
% \begin{commandline}
%	\begin{verbatim}
%		$ ls
%		
%		Applications	Desktop	...
%	\end{verbatim}
% \end{commandline}

\mdfdefinestyle{commandline}{
	leftmargin=10pt,
	rightmargin=10pt,
	innerleftmargin=15pt,
	middlelinecolor=black!50!white,
	middlelinewidth=2pt,
	frametitlerule=false,
	backgroundcolor=black!5!white,
	frametitle={Command Line},
	frametitlefont={\normalfont\sffamily\color{white}\hspace{-1em}},
	frametitlebackgroundcolor=black!50!white,
	nobreak,
}

% Define a custom environment for command-line snapshots
\newenvironment{commandline}{
	\medskip
	\begin{mdframed}[style=commandline]
}{
	\end{mdframed}
	\medskip
}

%----------------------------------------------------------------------------------------
%	FILE CONTENTS ENVIRONMENT
%----------------------------------------------------------------------------------------

% Usage:
% \begin{file}[optional filename, defaults to "File"]
%	File contents, for example, with a listings environment
% \end{file}

\mdfdefinestyle{file}{
	innertopmargin=1.6\baselineskip,
	innerbottommargin=0.8\baselineskip,
	topline=false, bottomline=false,
	leftline=false, rightline=false,
	leftmargin=2cm,
	rightmargin=2cm,
	singleextra={%
		\draw[fill=black!10!white](P)++(0,-1.2em)rectangle(P-|O);
		\node[anchor=north west]
		at(P-|O){\ttfamily\mdfilename};
		%
		\def\l{3em}
		\draw(O-|P)++(-\l,0)--++(\l,\l)--(P)--(P-|O)--(O)--cycle;
		\draw(O-|P)++(-\l,0)--++(0,\l)--++(\l,0);
	},
	nobreak,
}

% Define a custom environment for file contents
\newenvironment{file}[1][File]{ % Set the default filename to "File"
	\medskip
	\newcommand{\mdfilename}{#1}
	\begin{mdframed}[style=file]
}{
	\end{mdframed}
	\medskip
}

%----------------------------------------------------------------------------------------
%	NUMBERED QUESTIONS ENVIRONMENT
%----------------------------------------------------------------------------------------

% Usage:
% \begin{question}[optional title]
%	Question contents
% \end{question}

\mdfdefinestyle{question}{
	innertopmargin=1.2\baselineskip,
	innerbottommargin=0.8\baselineskip,
	roundcorner=5pt,
	nobreak,
	singleextra={%
		\draw(P-|O)node[xshift=1em,anchor=west,fill=white,draw,rounded corners=5pt]{%
		Question \theQuestion\questionTitle};
	},
}

\newcounter{Question} % Stores the current question number that gets iterated with each new question

% Define a custom environment for numbered questions
\newenvironment{question}[1][\unskip]{
	\bigskip
	\stepcounter{Question}
	\newcommand{\questionTitle}{~#1}
	\begin{mdframed}[style=question]
}{
	\end{mdframed}
	\medskip
}

%----------------------------------------------------------------------------------------
%	WARNING TEXT ENVIRONMENT
%----------------------------------------------------------------------------------------

% Usage:
% \begin{warn}[optional title, defaults to "Warning:"]
%	Contents
% \end{warn}

\mdfdefinestyle{warning}{
	topline=false, bottomline=false,
	leftline=false, rightline=false,
	nobreak,
	singleextra={%
		\draw(P-|O)++(-0.5em,0)node(tmp1){};
		\draw(P-|O)++(0.5em,0)node(tmp2){};
		\fill[black,rotate around={45:(P-|O)}](tmp1)rectangle(tmp2);
		\node at(P-|O){\color{white}\scriptsize\bf !};
		\draw[very thick](P-|O)++(0,-1em)--(O);%--(O-|P);
	}
}

% Define a custom environment for warning text
\newenvironment{warn}[1][Warning:]{ % Set the default warning to "Warning:"
	\medskip
	\begin{mdframed}[style=warning]
		\noindent{\textbf{#1}}
}{
	\end{mdframed}
}

%----------------------------------------------------------------------------------------
%	INFORMATION ENVIRONMENT
%----------------------------------------------------------------------------------------

% Usage:
% \begin{info}[optional title, defaults to "Info:"]
% 	contents
% 	\end{info}

\mdfdefinestyle{info}{%
	topline=false, bottomline=false,
	leftline=false, rightline=false,
	nobreak,
	singleextra={%
		\fill[black](P-|O)circle[radius=0.4em];
		\node at(P-|O){\color{white}\scriptsize\bf i};
		\draw[very thick](P-|O)++(0,-0.8em)--(O);%--(O-|P);
	}
}

% Define a custom environment for information
\newenvironment{info}[1][Info:]{ % Set the default title to "Info:"
	\medskip
	\begin{mdframed}[style=info]
		\noindent{\textbf{#1}}
}{
	\end{mdframed}
}
 % Include the file specifying the document structure and custom commands

%----------------------------------------------------------------------------------------
%	ASSIGNMENT INFORMATION
%----------------------------------------------------------------------------------------

\title{Tarefa III PCS5028 - Metodologia de Pesquisa Científica em Engenharia de Computação} % Title of the assignment

\author{Giovani Tavares\\ \texttt{giovanitavares@usp.br}} % Author name and email address

\date{Universidade de São Paulo --- July 26, 2026} % University, school and/or department name(s) and a date

%----------------------------------------------------------------------------------------

\begin{document}

\maketitle % Print the title

%----------------------------------------------------------------------------------------
%	INTRODUCTION
%----------------------------------------------------------------------------------------


\section*{Submeter breve relato sobre discussão em aula, a sua percepção, avaliando o que é preciso amadurecer/melhorar no seu projeto.}

Em especial, três aspectos discurtidos em aula me levaram a reavaliar alguns aspectos do meu projeto de pesquisa. 

O primeiro refere-se à metodologia de pesquisa, a qual deve refletir a sequência de passos necessária para se atingir o objetivo principal e, portanto, também contem a sequência de passos para refletir os objetivos secundários relacionados ao principal. Até então, a metodologia que vinha escrevendo em meus rascunhos estava mais relacionada à sequência de tarefas que venho realizando desde o início do meu mestrado, a qual não reflete \textit{somente} o caminho para atingir meu objetivo, apesar de também incluí-lo. Eu havia escrito na metodologia as etapas de definição do objetivo em si, as quais devem fazer parte das iterações de revisão bibliográfica discutidas em aula e não necessariamente precisam contar no exame de qualificação ou na dissertação final.

Sendo assim, retirei da metodologia todas as discussões sobre como eu cheguei ao objetivo de pesquisa e todos os experimentos que realizei para chegar nele. No lugar, deixei apenas as etapas que completei e ainda irei completar para atingí-lo. 

O segundo ponto que merece amadurecimento em meu projeto e exame de qualificação refere-se a justificativa de minha hipótese de pesquisa. Eu a considerava minha hipótese como sendo válida e digna de ser investigada pelo fato de que em minhas iterações de revisão bibliográfica eu não encontrei pesquisas que avaliassem o uso de ontologias e \textit{large language models} (LLMs)  no domínio específico que estou estudando. Por outro lado, as discussões em sala de aula me levaram a adicionar o \textit{porquê} esse uso pode fazer sentido, e isso se explica nas minhas próprias referências principais que aplicam as técnicas que estudo em outros domínios. Ou seja, melhorei a forma de justificar a minha hipótese explicitando pesquisas em que o uso de ontologias com LLMs se demonstrou bem sucedido em outros domínios. 

Por último, a atividade em grupo ao final da aula, em que tivemos que avaliar o projeto de pesquisa de outro grupo inciou uma discussão interessante sobre o que deveria ser considerado ou não uma limitação do trabalho avaliado. Discutimos se a falta de dados anotados para o trabalho era uma limitação ou apenas uma dificuldade, dado que esses dados poderiam ser obtidos, sendo tal obtenção parte da metodologia para se atingir o objetivo proposto. Não chegamos a um consenso no grupo, mas a discussão me fez refletir sobre esse aspecto em minha própria pesquisa, uma vez que eu tenha dados para realizá-la, entretanto o fato de que eles não são abundantes no formato que preciso me levou a reduzir o escopo daquilo que estou pesquisando.


\section*{Identificar as 3 principais referências para o seu trabalho, justificando}

As três mais importantes referências para meu trabalho são:

1. Sharma et al. (2025) — OG-RAG** `arXiv:2412.15235`
Demonstra que recuperar informação sobre uma ontologia formal pré-definida produz respostas mais corretas do que recuperar sobre texto plano ou sobre grafos extraídos automaticamente do corpus. Isso justifica diretamente a hipótese central deste trabalho: que responder perguntas de domínio sobre GEMs usando OntoGEM como estrutura de recuperação produz respostas de maior qualidade do que sistemas que ignoram a ontologia.

2. Edge et al. (2024) — GraphRAG** `arXiv:2404.16130`
Para afirmar que *ontologia* melhora as respostas, não basta comparar com texto plano — é preciso mostrar que a melhora vai além do que qualquer estrutura de grafo (mesmo sem schema formal) já proporcionaria. GraphRAG é o método que fornece essa estrutura sem ontologia, tornando possível isolar se o ganho vem da formalidade do schema ou simplesmente de haver um grafo.

3. King et al. (2020) — BiGG Models** `NAR 48:D402`
As perguntas de domínio sobre GEMs frequentemente exigem combinar dados estruturados (regras GPR, estequiometria) com texto de publicações científicas. A hipótese de que a ontologia degrada menos nessa combinação do que os baselines só é testável num corpus que tenha os dois lados — e BiGG é esse corpus, com modelos curados, API estável e referências às publicações originais.

\end{document}
